%   Harness 会影响 agent performance
%                       ↓
%   现有 harness optimization 通常通过专门的优化过程，
%   从任务轨迹和，对照明确的 reference 学习并改进 harness, 
%   用预定义好的 success rate 或者 end-state outcome 进行评估
%                       ↓
%   但许多真实任务是 open-ended的
%                       ↓
%   提出 PerfRace，评价一个 run 内持续产生的进展
%                       ↓
%   我们进一步研究 harness optimization 能不能作为一种在长程任务中
%   由 agent 控制的、递归的能力
%                       ↓
%   提出 BOXBOX (F1 赛车里面汽车进站维修的行为，我们用它来比喻 agent 可以一边跑一边“换轮子”)：
%   agent-triggered、in-place、runtime harness optimization

\begin{abstract}
Automated harness optimization improves agents by learning from trajectories and task feedback. Existing offline methods use reference answers to optimize harnesses for reuse on held-out tasks. However, many real-world tasks are open-ended, while long runs expose new difficulties that offline-optimized harnesses may not address. To study harness optimization in these settings, we introduce \perfrace{}, a benchmark for long-running performance optimization of software projects, and investigate whether agents can use task-specific experience to recursively improve their harnesses during an ongoing task. To this end, we introduce \boxbox{}, which enables agent-triggered harness revision and kernel-managed hot reloading while preserving execution context. We evaluate \boxbox{} on \perfrace{} under matched test-time budgets, comparing it with Frozen Harness, offline Harness Evolution, and Harness Scaling. \boxbox{} achieves higher speedup with more efficient use of task time and model cost, reaching a \speedupat{4\mathrm{h}} of \(1.70\times\), compared with \(1.54\times\) for the strongest baseline.
% Code is available at \url{https://anonymous.4open.science/r/BoxBox-1C35/}.
\end{abstract}
